\subsection{Pointed-associator rigidity for \texorpdfstring{$\sha{R8_3_65}$}{R8-3-65}}
\label{sec:r65-pointed-associator-next}
This refinement removes categorical branches of the connected-algebra
search, while leaving the fusion ring unresolved.

\begin{proposition}
In any complex categorification of $\sha{R8_3_65}$, the pointed subcategory
has nonzero class $k\in H^3(C_4,\C^\times)\cong\Z/4$.
Neither of the two rank-two based module actions can be categorified.
Either rank-four action requires $k=2$; if $k=2$, the subgroup algebra
$\one\oplus g^2$ has a Morita dual with the same Grothendieck ring
$\sha{R8_3_65}$ and pointed class $2$.
\end{proposition}
\begin{proof}
If $k=0$, the actual subgroup algebra for $H=\langle g^2\rangle$ exists.
Its four simple right modules arise from the four free right $H$-orbits.
Each has two left $H$-linearizations, since the possible projective
multiplier is trivializable by $H^2(C_2,\C^\times)=0$.
The bimodule dual has four invertibles forming
$\widehat H\times(C_4/H)\cong V_4$, and a free orbit of four simples of
dimension $d=(11+5\sqrt5)/2$, where $d^2=1+11d$.
This argument uses a subgroup algebra, without assuming a central or
Tannakian lift.

For a free noninvertible orbit over an abelian group $G$ of order four,
write $\rho g=\phi(g)\rho$ and
$\rho^2=u+(\sum_h n_hh)\rho$, with $\sum_hn_h=11$.
Associativity forces $\phi^2=1$ and invariance of $n_h$ under
$\{g\phi(g)^{-1}\}$. If $\phi\ne1$, this subgroup has order two and would
make the coefficient sum even. Thus $\phi=1$.
Duality and the odd coefficient sum permit a translate with $\rho^*=\rho$,
$\rho^2=1+N\rho$ and $N=N^*$.
For $V_4$ the complete possibilities are the $27$ unordered partitions of
$11$ into four parts. For a character parameter $m=\chi(N)$,
\[
 \operatorname{Norm}(f)=16(m^2+4),\qquad
 \operatorname{Norm}(\FPdim(\mathcal D))=2000.
\]
The two-FP-value Galois/pseudo-unitary reduction applies to this dual.
Integrality of its distinguished central dimension $D/f$, taking norms in
the quadratic compositum, implies $125/(m^2+4)\in\Z$.
Hence $|m|=1$ or $11$. Only the partitions
$(0,0,0,11)$, $(0,0,5,6)$ and $(2,3,3,3)$ remain. They are excluded,
respectively, by the general rank-two classification, the repository's
unrestricted rank-four dimension bound, and the exact identification with
the already excluded $\sha{R8_3_1453}$.
The accepted rank-four and $\sha{R8_3_1453}$ results are explicit inputs to
this deduction. Thus $k\ne0$.

In a rank-two action each invertible fixes every module simple, so its
scalar module associators would trivialize $k$. In a rank-four action the
stabilizers are $C_2$, so restriction of $k$ to $H^3(C_2,\C^\times)$ is zero.
The restriction map is reduction modulo two, forcing $k=2$.
Conversely, for $k=2$ the subgroup algebra exists. Its dual has an
order-four invertible group and the same free noninvertible orbit.
The complete $V_4$ comparison above and the corresponding $21$ selfdual
$C_4$ coefficient vectors leave only the ring $\sha{R8_3_65}$.
The rank-four Morita module on the dual side forces its pointed class to
be $2$ as well.
\end{proof}
The categorical cuts reduce the four quadratic rings' aggregate necessary
lists to $144$ algebra objects, $178$ induction factors and $30$ full-action
classes. For $\sha{R8_3_65}$ the remaining action ranks are $4,4,8,8$.
The conditional subgroup dual supplies neither an associator for this ring
nor an equivalence between the dual and the original category.
The complete finite coefficient domains, theorem dependencies and source
hashes are recorded in
\path{evidence/rank8_next_20261021/quadratic/POINTED_COCYCLE.md} and
\path{pointed_dual.py} in that directory, with an independent categorical
review in
\path{evidence/rank8_next_20261021/semion_fermion/construction/R65_ADDENDUM_REVIEW.md}.
